\chapter{Búsquedas}

\section{Búsqueda binaria}
Siempre si vas a resolver algo con lower\_bound aunque lo diga o no lo diga el problema SIEMPRE pero SIEMPRE ordenalo, me comi varios wa porque confiaba que siempre venia ordenado. \\ 
La función \texttt{lower\_bound} es una herramienta fundamental en la biblioteca estándar de C++ para realizar búsquedas eficientes en contenedores ordenados.

\subsection{Función \texttt{lower\_bound}}
\begin{itemize}
	\item \textbf{Complejidad}: $O(\log n)$ para contenedores de acceso aleatorio como \texttt{vector}
	\item \textbf{Comportamiento}: Devuelve un iterador al primer elemento que \textbf{no es menor} que el valor dado (primer elemento $\geq$ al valor buscado)
	\item \textbf{Sintaxis}: \texttt{lower\_bound(iterador\_inicio, iterador\_fin, valor)}
\end{itemize}

\section{Problema de ejemplo}

\subsection{Descripción}
El aula se representa como una línea unidimensional con celdas desde $1$ hasta $n$. Inicialmente hay $m$ profesores y David en posiciones distintas. En cada movimiento:
\begin{enumerate}
	\item David se mueve a una celda adyacente o se queda
	\item Todos los profesores se mueven simultáneamente a una celda adyacente o se quedan
\end{enumerate}

David es atrapado cuando algún profesor está en la misma celda. Todos actúan óptimamente: David maximiza el número de movimientos hasta ser atrapado, mientras los profesores lo minimizan.

\subsection{Solución}
Para cada consulta con la posición $a_i$ de David:
\begin{enumerate}
	\item Encontrar los dos profesores más cercanos (izquierda y derecha)
	\item Calcular la mitad de la distancia entre ellos
	\item David se moverá al punto medio para maximizar el tiempo
\end{enumerate}

\subsection{Código de Solución}
\begin{lstlisting}[style=cppstyle, caption={Solución del problema}]

	void solve() {
		ll n; cin >> n;
		ll m, q; cin >> m >> q;
		vi b(m);
		fore(i, 0, m) cin >> b[i];
		vi a(q);
		fore(i, 0, q) cin >> a[i];
		
		sort(all(b));
		
		fore(i, 0, q) {
			ll izq = 0, der = 0;
			ll res = 0;
			auto it = lower_bound(all(b), a[i]);
			
			if (it == b.begin()) {
				res = (*it) - 1;
			} else if (it == b.end()) {
				--it;
				res = n - (*it);
			} else {
				der = *it;
				--it;
				izq = *it;
				ll dist_der = der - a[i];
				ll dist_izq = a[i] - izq;
				res = (dist_der + dist_izq) / 2;
			}
			pri(res);
		}
	}
\end{lstlisting}

\section{Uso de \texttt{lower\_bound} para obtener índices}

\subsection{Obtención del índice}
Cuando se utiliza \texttt{lower\_bound} sobre un contenedor, podemos obtener la posición (índice) del elemento encontrado mediante aritmética de punteros:

\begin{lstlisting}[style=cppstyle, caption={Obtención del índice con lower\_bound}]
	void solve() {
		ll n; cin>>n; ll k; cin>>k;
		vi a(n); fore(i, 0, n) cin>>a[i];
		
		while(k--) {
			ll val; cin >> val;
			auto it = lower_bound(all(a), val);
			ll res = it - a.begin();  // Obtenemos el indice
			pri(res + 1);  // +1 para indice base 1
		}
	}
\end{lstlisting}

\subsection{Problema}
Tenemos un edificio de $n$ pisos, cada piso $i$ con $a_i$ departamentos. Las cartas están numeradas linealmente desde $1$ hasta la suma total de departamentos.  
Dado un número de carta $b_j$, hay que determinar:
\begin{itemize}
	\item El \textbf{piso} donde va la carta.
	\item El \textbf{departamento dentro del piso}.
\end{itemize}

\textbf{Ejemplo:} $n=2$, $a=[3,5]$, carta $b=7$ → Piso $2$, departamento $4$.

Idea de la solución
1. Construimos un vector de \textbf{prefix sums}:
\[
\text{pref}[0] = a_0, \quad
\text{pref}[i] = \text{pref}[i-1] + a_i \quad \text{para } i \ge 1
\]
Cada elemento indica la cantidad total de departamentos hasta ese piso.  

2. Para cada carta $b_j$ usamos \textbf{búsqueda binaria} con \texttt{lower\_bound}:
\[
\text{piso} = \text{lower\_bound}(\text{pref.begin(), pref.end(), } b_j) - \text{pref.begin()}
\]

3. Calculamos el departamento dentro del piso:
\[
\text{depto} = b_j - (\text{piso}>0 ? \text{pref[piso-1]} : 0)
\]

---

Código en C++
\begin{verbatim}

	void solve() {
		ll n, m;
		cin >> n >> m;
		vector<ll> a(n), b(m);
		for (int i = 0; i < n; i++) cin >> a[i];
		for (int i = 0; i < m; i++) cin >> b[i];
		
		// Construir prefix sums
		vector<ll> pref(n);
		pref[0] = a[0];
		for (int i = 1; i < n; i++) pref[i] = pref[i-1] + a[i];
		
		for (int i = 0; i < m; i++) {
			// Encontrar piso usando búsqueda binaria
			int piso = lower_bound(pref.begin(), pref.end(), b[i]) - pref.begin();
			// Departamento dentro del piso
			ll depto = b[i] - (piso > 0 ? pref[piso-1] : 0);
			cout << piso + 1 << " " << depto << "\n";
		}
	}
	
	
\end{verbatim}



Ejemplo paso a paso
\[
a = [3,5,2] \implies pref = [3,8,10], \quad b = [1,4,9]
\]

\begin{tabular}{c|c|c|c}
	Carta $b_j$ & Piso (0-based) & Piso (1-based) & Departamento \\
	\hline
	1 & 0 & 1 & 1 \\
	4 & 1 & 2 & 1 \\
	9 & 2 & 3 & 1
\end{tabular}

Complejidad
\begin{itemize}
	\item Construir prefix sums: $O(n)$
	\item Cada carta: $O(\log n)$ con \texttt{lower\_bound}
	\item Total: $O(n + m \log n)$
\end{itemize}

Tips
\begin{itemize}
	\item `lower\_bound` devuelve iterador, se resta `pref.begin()` para obtener índice.
	\item Esta técnica combina \textbf{prefix sums} con \textbf{búsqueda binaria} para ubicar rápidamente rangos acumulados.
\end{itemize}

\section{Búsqueda binaria en la respuesta}

\subsection{Caso general}
\begin{lstlisting}[style=cppstyle, caption={Busqueda binaria en la respuesta}]
	bool can(int mid) {
		// Esta funcion debe responder si con mid como respuesta
		// es posible cumplir con la condicion del problema
		
		// Implementar la logica del chequeo
		
		return true;
	}
	
	int main(){
		// Busqueda binaria
		int low = 0, high = 1e9; // Ajustar limites segun el problema
		int ans = -1;
		
		while (low <= high) {
			int mid = (low + high) / 2;
			
			// si queremos el minimo que cumple
			if (check(mid)) {
				ans = mid;        
				high = mid - 1;
			} else {
				low = mid + 1;
			}
			
			// Si queres el maximo que cumple:
			if (check(mid)) {
				ans = mid;
				low = mid + 1;
			} else {
				high = mid - 1;
			}
		}
		
		cout << ans << endl;
	}
	
\end{lstlisting}

\subsection{Ejemplo 1: Problema de Sumas}
\begin{lstlisting}[style=cppstyle, caption={Problema de sumas con búsqueda binaria}]
	ll n;
	bool can(ll x, vll &a) {
		ll cont = 0; // Sup. disponibles
		fore(i, 0, n) {
			if (x < a[i]) return false; 
			cont += x - a[i];
		}
		return cont >= x;
	}
	
	void solve() {
		cin >> n; 
		vll a(n); 
		fore(i, 0, n) cin >> a[i]; 
		
		ll l = 0, r = 1e14, ans = 0;
		while (l <= r) {
			ll mid = (l + r) >> 1; 
			if (can(mid, a)) {
				ans = mid; 
				r = mid - 1;
			} else {
				l = mid + 1;
			}
		} 
		pri(ans);
	}
\end{lstlisting}

\subsection{Ejemplo 2: Problema de Porcentajes}
\begin{lstlisting}[style=cppstyle, caption={Problema de porcentajes con tickets}]
	ll k, x, y, a, b, n;
	vll pr;
	
	bool can(ll st) {
		vll p(st);
		for (int i = a - 1; i < st; i += a) p[i] += x;
		for (int i = b - 1; i < st; i += b) p[i] += y;
		sort(all(p), greater<ll>());
		ll cur = 0;
		fore(i, 0, st) cur += p[i] * pr[i] / 100;
		return cur >= k;
	}
	
	void solve() {
		pr.clear();
		cin >> n;
		fore(i, 0, n) {
			ll ai;
			cin >> ai;
			pr.pb(ai);
		}
		sort(all(pr), greater<int>()); 
		cin >> x >> a >> y >> b >> k;
		
		ll l = 0, h = n, ans = -1;
		while (l <= h) {
			ll mid = l + (h - l) / 2;
			if (can(mid)) {
				ans = mid;
				h = mid - 1;
			} else {
				l = mid + 1;
			}
		}
		pri(ans);
	}
\end{lstlisting}

\subsection{Ejemplo 3: Problema de Recursos}
\begin{lstlisting}[style=cppstyle, caption={Problema de gestión de recursos}]
	ll k, n; 
	vi a, b; 
	
	bool can(ll mid) {
		ll ka = k;
		fore(i, 0, n) {
			ll totact = a[i] * mid; 
			if (totact > b[i]) {
				ll rest = totact - b[i];
				if (rest <= ka) {
					ka -= rest;
				} else {
					return false;
				}
			}
		}
		return true;
	}
	
	void solve() {
		cin >> n >> k; 
		a.resize(n);
		b.resize(n);
		fore(i, 0, n) cin >> a[i]; 
		fore(i, 0, n) cin >> b[i]; 
		
		ll low = 0, high = 2e9 + 10;
		ll res = -1;
		
		while (low <= high) {
			ll mid = (low + high) / 2;
			if (can(mid)) {
				res = mid;
				low = mid + 1; // Buscamos el MAXIMO
			} else {
				high = mid - 1;
			}
		}
		pri(res);
	}
\end{lstlisting}

\subsection{Consideraciones Importantes}
\begin{itemize}
	\item \textbf{Límites}: Definir correctamente \texttt{l} y \texttt{r} según el problema
	\item \textbf{Función can}: Debe ser eficiente (generalmente $O(n)$ o $O(n \log n)$)
	\item \textbf{Overflow}: Usar \texttt{mid = l + (r - l) / 2} para evitar overflow
	\item \textbf{Tipo de búsqueda}: Decidir si se busca el mínimo o máximo que cumple
\end{itemize}

\subsection{Complejidad}
La complejidad total es $O(T \times \log(\text{rango}) \times \text{complejidad de can})$, donde:
\begin{itemize}
	\item $T$: número de casos de prueba
	\item \texttt{rango}: diferencia entre \texttt{r} y \texttt{l}
	\item \texttt{can}: complejidad de la función de verificación
\end{itemize}


\section{Comparadores personalizados}

A veces necesitás ordenar de una forma específica, o hacer búsquedas binarias con criterios distintos al \texttt{<} por defecto. Para eso existen los \textbf{comparadores}.

Un comparador es una función (o lambda) que recibe dos elementos y devuelve \texttt{true} si el primero \textbf{va antes} que el segundo en el orden deseado.

\subsection{Comparadores en \texttt{sort}}

\subsubsection{Orden inverso con \texorpdfstring{\texttt{greater\textless T\textgreater}}{greater(T)}}

\begin{lstlisting}[style=cppstyle, caption={Ordenar de mayor a menor}]
    vector<int> v = {3, 1, 4, 1, 5, 9};
    sort(all(v), greater<int>());
    // v = {9, 5, 4, 3, 1, 1}
\end{lstlisting}

\subsubsection{Comparador con lambda}

\begin{lstlisting}[style=cppstyle, caption={Ordenar con lambda personalizada}]
    // Ordenar por segundo elemento del par, si empatan por primero
    vector<pair<int,int>> v = {{1,3},{2,1},{3,2}};
    sort(all(v), [](auto& a, auto& b){
        if (a.second != b.second) return a.second < b.second;
        return a.first < b.first;
    });
    // v = {{2,1},{3,2},{1,3}}
\end{lstlisting}

\subsubsection{Comparador como struct (necesario para \texttt{set} y \texttt{map})}

Para \texttt{set} y \texttt{map}, las lambdas no sirven como comparador de tipo porque no tienen un tipo deducible en tiempo de compilación. Hay que usar un \textbf{struct con \texttt{operator()}}:

\begin{lstlisting}[style=cppstyle, caption={Struct comparador para set}]
    struct Cmp {
        bool operator()(const pair<int,int>& a, const pair<int,int>& b) const {
            // Ordena por suma, si empatan por primer elemento
            if (a.first + a.second != b.first + b.second)
                return a.first + a.second < b.first + b.second;
            return a.first < b.first;
        }
    };

    set<pair<int,int>, Cmp> s;
    s.insert({2, 1});  // suma 3  <- primero
    s.insert({0, 4});  // suma 4
    s.insert({1, 3});  // suma 4, primer elem 1 <- segundo
\end{lstlisting}

\subsection{\texttt{lower\_bound} y \texttt{upper\_bound} con comparador}

\texttt{lower\_bound} y \texttt{upper\_bound} aceptan un comparador como cuarto argumento.
\textbf{Importante:} el vector debe estar ordenado con \textbf{el mismo comparador}.

\begin{lstlisting}[style=cppstyle, caption={lower\_bound con greater (vector de mayor a menor)}]
    vector<int> v = {9, 5, 4, 3, 1, 1};
    // Vector ordenado de mayor a menor

    // lower_bound con greater: primer elemento no mayor que val
    // Es decir: primer elemento <= val
    auto it = lower_bound(all(v), 4, greater<int>());
    cout << *it << "\n"; // 4

    // upper_bound con greater: primer elemento estrictamente menor que val
    it = upper_bound(all(v), 4, greater<int>());
    cout << *it << "\n"; // 3
\end{lstlisting}

\subsubsection{Equivalencias útiles}

\begin{center}
\begin{tabular}{l|l|l}
    \textbf{Quiero encontrar...} & \textbf{Orden asc. (default)} & \textbf{Orden desc. (\texttt{greater})} \\
    \hline
    Primer elemento $\geq x$ & \texttt{lower\_bound(..., x)} & \texttt{upper\_bound(..., x, greater<T>())} \\
    Primer elemento $> x$    & \texttt{upper\_bound(..., x)} & \texttt{lower\_bound(..., x, greater<T>())} \\
    Último $\leq x$ & \texttt{prev(upper\_bound(..., x))} & \texttt{lower\_bound(..., x, greater<T>())} \\
    Último $< x$    & \texttt{prev(lower\_bound(..., x))} & \texttt{upper\_bound(..., x, greater<T>())} \\
\end{tabular}
\end{center}

\subsection{\texttt{lower\_bound} en \texttt{set} y \texttt{map}}

Para \texttt{set} y \texttt{map} \textbf{no usar la versión global} de \texttt{lower\_bound} sino el \textbf{método del contenedor}: \texttt{s.lower\_bound(x)}. La versión global no conoce el comparador interno del set.

\begin{lstlisting}[style=cppstyle, caption={lower\_bound en set normal y set invertido}]
    set<int> s = {1, 3, 5, 7, 9};

    auto it = s.lower_bound(4); // primer elemento >= 4
    cout << *it << "\n"; // 5

    it = s.upper_bound(4);      // primer elemento > 4
    cout << *it << "\n"; // 5

    // Set en orden inverso
    set<int, greater<int>> sr = {1, 3, 5, 7, 9};
    // sr internamente: 9 7 5 3 1

    // Con greater, lower_bound busca el primer elemento
    // "no menor" segun greater, es decir: primer elemento <= val
    auto it2 = sr.lower_bound(6);
    cout << *it2 << "\n"; // 5

    // upper_bound: primer elemento estrictamente "menor" segun greater,
    // es decir: primer elemento < val
    it2 = sr.upper_bound(6);
    cout << *it2 << "\n"; // 5 (6 no estaba, mismo resultado)

    it2 = sr.upper_bound(7);
    cout << *it2 << "\n"; // 5 (primer elemento < 7 en orden desc)
\end{lstlisting}

\subsection{Comparadores para pares}

Por defecto los pares se comparan lexicográficamente. Ejemplo de cómo modificarlo:

\begin{lstlisting}[style=cppstyle, caption={Ordenar pares por segundo elemento}]
    vector<pair<int,int>> v = {{3,1},{1,4},{2,2}};

    // Ordenar por segundo elemento (ascendente)
    sort(all(v), [](auto& a, auto& b){ return a.second < b.second; });
    // v = {{3,1},{2,2},{1,4}}

    // Ordenar por segundo elemento (descendente), si empatan
    // por primer elemento descendente
    sort(all(v), [](auto& a, auto& b){
        if (a.second != b.second) return a.second > b.second;
        return a.first > b.first;
    });
    // v = {{1,4},{2,2},{3,1}}
\end{lstlisting}

\subsection{Trampas comunes}

\begin{itemize}
    \item \textbf{Global \texttt{lower\_bound} en set}: usar siempre \texttt{s.lower\_bound(x)}, nunca \texttt{lower\_bound(s.begin(), s.end(), x)} con un set con comparador custom.
    \item \textbf{Lambda en set}: las lambdas no se pueden usar como tipo de comparador en \texttt{set}/\texttt{map} directamente. Usar un struct con \texttt{operator()}.
    \item \textbf{Comparador inconsistente}: si ordenás con un comparador y después llamás a \texttt{lower\_bound} con otro, el resultado es basura.
    \item \textbf{Strict weak ordering}: el comparador \textbf{debe} ser un orden estricto débil: irreflexivo, asimétrico y transitivo. \texttt{<=} como comparador \textbf{no} cumple esto y da comportamiento indefinido.
\end{itemize}

